\documentclass[10pt]{article}

\usepackage[utf8]{inputenc}

\usepackage[T1]{fontenc}

\usepackage{amsmath}

\usepackage{amsfonts}

\usepackage{amssymb}

\usepackage[version=4]{mhchem}

\usepackage{stmaryrd}

\usepackage{bbold}

\usepackage{graphicx}

\usepackage[export]{adjustbox}

\graphicspath{ {./images/} }

\usepackage{mathrsfs}



\begin{document}

ПРОЕКТИВНАЯ ПРЯМАЯ - проективное пространство размерности 1; П. П., рассматриваемая как самостоятельный объект, является замкнутым одномерным многообразием. П. П. является своеобразным проективным пространством - на ней нет интересных отношений инцидентности, как у проективных пространств большей размерности. Единственным инвариантом П. П. служит число ее точек. П. п. наз. н е п р е р ы в н о й, д и с к р е т н о й или к о н е ч н ой, если она инцидентна со множеством точек мощности континуума, счетным или конечным соответственно.



П. П. наз. у п о р я д о ч е н н о й, если на ней задано отношение разделения двух пар различных точек. Предполагается, что разделение не зависит от порядка пар и порядка точек в парах и любая четверка различных точек разбиваөтся на две взаимно разделяющиеся пары единственным образом, а также принимается аксиома расположения, связывающая пять различных точек (см., напр., [1]). Упорядочение П. п. над полем $\mathbb{R}$ связано с упорядоченностью этого поля, а именно: пара точек $\{A, B\}$ разделяет пару $\{C, D\}$, если двойное отношение $(A, B ; C, D)$ отрицательно, и не разделяет, если $(A, B ; C, D)$ положительно. Конечную П. п. $P G(1, q)$ над Галуа полем нечетного порядка $q$ можно упорядочить аналогично вегцественной П. П. Полагают (см. $[4])$, что пара точек $\{A, B\}$ разделяет пару $\{C, D\}$ тогда и только тогда, когда $(A, B ; C, D)$ - квадратичный вычет поля Галуа $G F(q)$.



П. П. приобретает определенное геометрич. строение, если она вложена в проективное пространство большей размерности; так, напр., П. п. однозначно определяется двумя различными точками, а аналитич. определение П. П. как множества классов эквивалентности пар әлементов тела $k$, не равных одновременно нулю, по существу эквивалентно вложению П. П. в проективное пространство $P_{n}(k), k \geqslant 2$. Если $P_{1}(k)$ является П. П. над полем $k$, то группа автоморфизмов П. п. Aut $P_{1}(k)$ может быть представлена на точках $P_{1}(k)$ в параметрич. форме как множество отображений



\[

k \longrightarrow \frac{k^{\alpha} a+b}{k^{\alpha}{ }_{c+d}}, a, b, c, d \in k, a d-b c \neq 0, \alpha \in \text { Aut } k .

\]



Группа алгебраич. автоморфизмов действительной П. П. изоморфна группе перемещений действительной плоскости Лобачевского, а порядок группы Aut $P G\left(1, p^{h}\right)$ равен $h\left(p^{3 h}-p^{h}\right)$.



На П. П. можно построить другие геометрии. Так, напр., плоскость Мёбиуса порядка $p$ допускает интерпретацию на П. п. $P G\left(1, p^{2}\right)$ (см. [5]). Другой традиционной геометрич. конструкцией является изображение проективного пространства $P_{n}(k)$ на П. П. $P_{1}(k)$ (см. [2]), при к-ром точки из $P_{n}(k)$ изображаются набором $n$ точек П. П. $P_{1}(k)$ (здесь $k$ - алгебраически замкнутое поле).



Лит.: [1] $\Gamma$ л а г о л е в H. А., Проективная геометрия, М., 1963; [2] III а ф а р ев ич И. Р., Основы алгебраической



\includegraphics[max width=\textwidth, center]{2023_07_08_41c30abfddbc93a3e75eg-1(1)}

Projective planes, N. Y., 1973; [4] K u s t a a n h e i m o P.,

"Comment. phys-math.", 1957, v. 20, № 8, [5] V e b l e n O.,



\begin{center}

\includegraphics[max width=\textwidth]{2023_07_08_41c30abfddbc93a3e75eg-1}

\end{center}



ІРОЕКТИВНАЯ СВЯЗНОСТЬ- - диференциально-геометрическая структура на гладком многообразии $M$; специальный вид связности на многообразии, когда приклеенное к $M$ гладкое расслоенное пространство $E$ имеет своим типовым слоем проективное пространство $P_{n}$ размерности $n=\operatorname{dim} M$. Структурой такого $E$ к каждой точке $x \in M$ присоединяется әкземпляр проективного пространства $\left(P_{n}\right)_{x}$, к-рый отождествляется (с точностью до гомологий с инвариантной связкой прямых в точке $x$ ) с касательным центроаффинным пространством $T_{x}(M)$, дополненным бесконечно удаленной гиперплоскостью. П. с., как связность в таком $E$, предусматривает сопоставление каждой гладкой кривой $\mathscr{L} \in M$ с началом $x_{0}$ и каждой ее точке $x_{t}$ проективного отображения $\left(P_{n}\right)_{x_{t}} \rightarrow\left(P_{n}\right)_{x_{0}}$ так, что удовлетворяется следующее условие. Пусть $M$ покрыто координатными областями, в к-рых фиксировано гладкое поле репера $\left(P_{n}\right)_{x}$, у к-рого вершина, определяемая вектором $e_{0}$, совпадает с $x$. (Репер в $P_{n}$ определяется классом эквивалентности базисов в векторном пространстве $V_{n+1}$, если эквивалентными считаются те $\left\{e_{\alpha}\right\}$ и $\left\{e_{\alpha}^{\prime}\right\}, \alpha=0,1$, $\ldots, n$, у к-рых $e_{\alpha}^{\prime}=\lambda e_{\alpha}, \lambda \neq 0$.) Тогда при $t \rightarrow 0$ отображение семейства должно стремиться к тождественному отображению, причем главная часть его отклонения от последнего должна определяться относительно поля реперов в нек-рой окрестности точки $x_{0}$ матрицей линейных дифференциальных форм



\[

\begin{gathered}

\omega_{\alpha}^{\beta}=\Gamma_{\beta i}^{\alpha} d x^{i}, \operatorname{det}\left\|\Gamma_{0 i}^{j}\right\| \neq 0, \\

\alpha, \beta=0,1, \ldots, n ; i, j=1, \ldots, n,

\end{gathered}

\]



общей для всех $L$. Другими словами, образ репера в точке $x_{t}$ при отображении $\left(P_{n}\right)_{x_{t}} \rightarrow\left(P_{n}\right)_{x_{0}}$ должен быть определен векторами



\[

e_{\beta}\left[\delta_{\alpha}^{\beta}+\omega_{\alpha}^{\beta}(X) t+\varepsilon_{\alpha}^{\beta}(t)\right],

\]



где $X_{-}$- касательный вектор к $L$ в точке $x_{0}$ и $\lim _{t \rightarrow 0} \frac{\varepsilon_{\alpha}^{\beta}(t)}{t}=0$. Возможность перехода к әквивалентным базисам приводит к тому, что среди форм (1) существенны только



\[

\omega_{0}^{i}, \theta_{i}^{j}=\omega_{i}^{i}-\delta_{i}^{i} \omega_{0}^{0}, \omega_{i}^{0}

\]



При преобразовании репера поля в произвольной точке $x \in M$ согласно формулам $e_{\alpha}^{\prime}=A_{\alpha}^{\beta} e_{\beta}, \quad e_{\beta}=A_{\beta}^{\alpha^{\prime}} e_{\alpha^{\prime}}$, где $A_{0}^{i}=A_{0}^{i^{\prime}}=0$, т. е. при переходе к произвольному элементу главного расслоенного пространства П реперов в пространствах $\left(P_{n}\right)_{x}$, формы (1) заменяются следующшми 1-формами на П:



2-формы



\[

\omega_{\alpha^{\prime}}^{\beta^{\prime}}=A_{\gamma}^{\beta^{\prime}} d A_{\alpha^{\prime}}^{\gamma}+A_{\alpha^{\prime}}^{\gamma} A_{\delta}^{\beta^{\prime} \omega_{\gamma}^{\delta}}

\]



\[

\Omega_{\alpha^{\prime}}^{\beta^{\prime}}=d \omega_{\alpha^{\prime}}^{\beta^{\prime}}+\omega_{\gamma^{\prime}}^{\beta^{\prime}} \wedge \omega_{\alpha^{\prime}}^{\gamma^{\prime}}

\]



являются полубазовыми, т. е. линейными комбинациями $\omega_{0}^{k} \wedge \omega_{0}^{l}$, и тензорными, т. е. при преобразовании репера матрицами $A_{\alpha}^{\gamma}$ имеют место формулы



\[

\Omega_{\alpha}^{\beta^{\prime}}=A_{\alpha}^{\nu}, A_{\delta}^{\beta \prime} \Omega_{\gamma}^{\delta},

\]



где $\Omega_{\alpha}^{\beta^{\prime}}$, составлены из (3) аналогично (4). Для существенных форм (2) имеют место структурные уравнения П. с. (где для простоты опущены штрихи):



\[

\left.\begin{array}{c}

d \omega_{0}^{i}-\theta_{j}^{i} \wedge \omega_{0}^{j}=\Omega_{0}^{i} \\

d \theta_{i}^{i}+\theta_{k}^{j} \wedge \theta_{i}^{k}+\omega_{0}^{k} \wedge\left(\delta_{i}^{j} \omega_{k}^{0}-i \delta_{k}^{j} \omega_{i}^{0}\right)=\Theta_{i}^{j}, \\

d \omega_{i}^{0}+\omega_{k}^{0} \wedge \theta_{i}^{k}=\Omega_{i}^{0},

\end{array}\right\}

\]



где $\Theta_{i}^{j}=\Omega_{i}^{j}-\delta_{i}^{j} \Omega_{0}^{0}$. Здесь правые части полубазовы; они составляют систему форм кручения-кривизны П. с.



Равенство $\Omega_{0}^{i}=0$ имеет инвариантный смысл. В этом случае говорят о П. с. нулевого кручения; для нее $\Theta_{\imath^{\prime}}^{i,}=\Theta_{i}^{i}$. Инвариантные тождества



\[

\begin{gathered}

\Omega_{0}^{i} \equiv 0, \quad \Theta_{i}^{i} \equiv 0, \\

K_{i k j}^{j} \equiv 0, \Theta_{i}^{i}=\frac{1}{2} K_{i k l}^{j} \omega_{0}^{k} \wedge \omega_{0}^{l},

\end{gathered}

\]



выделяют специальный класс П. с., называемых (по Картану) н о р м а л ь н ы м и П. с.



Формы (1) определяют П. с. на $M$ однозначно: образ при отображении $\left(P_{n}\right)_{x_{t}} \rightarrow\left(P_{n}\right)_{x_{0}}$ репера в точке $x_{t}$





\end{document}